\chapteropen{10}

\begin{cele}
\item policzyć, ile godzin doby zostaje, gdy blok snu trwa 12 godzin, a dziecko, podróż albo choroba zabierają kolejne bloki;
\item odróżnić awarię rytmu od porażki planu i wiedzieć, co z każdej z nich wynika dla następnego dnia;
\item stosować zasadę powrotu bez nadrabiania i nie zamieniać jednej złej nocy w nowy harmonogram;
\item przeprowadzić miesięczny przegląd na pięciu pytaniach i podjąć z niego jedną zmianę;
\item ocenić, czy 12 godzin snu jest nadal wyborem, i przejść na 9 godzin bez poczucia, że przegrałeś.
\end{cele}

\section{Plan, który się nie domyka}

Rozdziały od pierwszego do dziewiątego zakładały dobę, którą da się ułożyć. Ten rozdział zakłada coś mniej wygodnego: zwykle się nie ułoży. Plan dobowy jest modelem, a każdy model ma założenia. Ten zakłada, że nikt nie budzi się w nocy bez powodu, że autobus przyjeżdża o czasie, że w lodówce jest woda i że w kalendarzu nie ma niczego, czego nie przewidziano. Każde z tych założeń kiedyś przestaje być prawdziwe, zwykle w tym samym tygodniu.

Najuczciwiej zacząć od dzieci. Przy małym dziecku blok 12 godzin snu w jednej części nie jest realny, a udawanie, że jest, zamienia plan w źródło poczucia winy. Niemowlęta budzą się w nocy, a ich sen ma charakter przerywany: \verify{u niemowląt sen przez pierwsze miesiące składa się z krótkich cykli i wymaga nocnych karmień lub uspokajania}. Kto wpisuje do tabeli dwanaście godzin w jednym bloku, dostaje tabelę, której nie da się wypełnić, i co wieczór ogląda jej porażkę.

Zmiana polega na tym, że zmienia się jednostka. Zamiast budżetu dobowego liczysz budżet tygodniowy. Dwanaście godzin dziennie daje $12 \cdot 7 = 84$ godziny tygodniowo. Tyle samo może wyjść z jednej długiej nocy, z dwóch krótszych nocy i z drzemki w ciągu dnia. Liczy się suma, a nie kształt. Plan pozwala wtedy na noc gorszą od pozostałych, o ile tydzień jako całość nie spada poniżej celu.

\begin{przyklad}[Tydzień z niemowlęciem]
Przykład ilustracyjny, nie czyjeś dane. Cel na trzy doby to $3 \cdot 12 = 36$ godzin snu. Noc pierwsza: 5 godzin ciągłego snu. Noc druga: 6 godzin z dwiema pobudkami. Drzemka w ciągu dnia: 1 godzina. Noc trzecia: 7 godzin. Razem: $5 + 6 + 1 + 7 = 19$ godzin. Brakuje $36 - 19 = 17$ godzin.

Plan tego nie zamknie i nie powinien udawać, że zamyka. Sama liczba 17 godzin nie mówi jednak, co robić dalej. Mówi tylko, że cel na ten tydzień był ustawiony źle. Następny tydzień zaczyna się od nowego celu, a nie od odrabiania straty.
\end{przyklad}

\begin{uwaga}[Kiedy to nie jest kwestia planu]
Jeśli deficyt utrzymuje się tygodniami i towarzyszy mu senność w ciągu dnia, wyraźne zmiany nastroju albo chwile zaśnięcia w sytuacjach, w których nie wolno zasypiać, to nie jest problem do rozwiązania tabelą. Wtedy właściwy jest kontakt z lekarzem albo z pielęgniarką środowiskową. Żaden harmonogram tego nie zastąpi.
\end{uwaga}

\section{Podróż i zmiana stref}

Podróż zabiera trzy rzeczy naraz: godziny, bo zmienia się strefa czasowa, miejsce, bo kuchnia, łóżko i lodówka są cudze, oraz czas na pakowanie i dojazd, który w tabeli zwykle nie istnieje. Każda z nich wymaga osobnej decyzji.

Zmiana strefy przesuwa zegar, a nie organizm. Często powtarzana zasada mówi o jednym dniu na każdą przesuniętą strefę. \verify{tempo przystosowania rytmu dobowego po zmianie stref zależy od kierunku podróży i jest zwykle opisywane jako około dnia na strefę, przy dużym rozrzucie między osobami}. Nie traktuj tego jako prawa. Traktuj jako punkt wyjścia do obserwacji własnego ciała.

Są dwie strategie i żadna nie jest dobra w ogóle, obie są dobre dla konkretnej podróży. Pierwsza: trzymasz zegar domowy, czyli jesz i śpisz według godziny z domu. Blok snu wypada wtedy w środku nocy miejscowej, wymaga zaciemnionego pokoju i tolerancji na dzienną senność, ale nie wymaga żadnej korekty w tabeli po powrocie. Druga: przestawiasz się od razu na godziny miejscowe. Lepsza, jeśli masz na miejscu wieczorne zobowiązania. Gorsza, jeśli wracasz po tygodniu i zaczynasz wszystko od nowa.

Prosty przelicznik: przy przelocie sześć godzin na wschód, gdy w miejscu jest 22:00, w domu jest 16:00. Zanim zaplanujesz pierwszy wieczór, przelicz godziny w obu strefach. To jedyna czynność, której nie da się pominąć, a najczęściej się ją pomija.

Kuchnia wymaga osobnej uwagi. Przepis na chleb drożdżowy na zimno ma 12 godzin fermentacji w lodówce, a masa ciasta z tabeli to \val{bread.1.mass} g. Jeśli miałbyś zacząć ciasto w lodówce wynajętego mieszkania, pytaniem nie jest, czy zdążysz, tylko czy lodówka będzie wolna i czy nie zajmie jej ktoś inny. W podróży nie piecz. Zostaw zakwas w domu, w lodówce, i wróć do wypieku po powrocie. Na czas wyjazdu wystarczy chleb kupiony na miejscu.

\section{Choroba, deadline i awaria rytmu}

Choroba zmienia cel, a nie tylko tempo. Przy gorączce nie planujesz bloków, tylko odpoczynek. Organizm sam wydłuża sen i nie jest to porażka planu. \verify{w czasie infekcji organizm zwykle potrzebuje dłuższego odpoczynku, a sen jest jednym z mechanizmów powrotu do zdrowia}. Budżet 12 godzin wraca dopiero wtedy, gdy wraca zdrowie. Próba utrzymania go w czasie choroby zwykle daje tylko dodatkowe zmęczenie.

Deadline działa inaczej, bo jest krótki i zwykle sam się zamyka. Trzy noce po 9 godzin dają $3 \cdot (12 - 9) = 9$ godzin deficytu. To jest koszt, który da się zapisać i policzyć. Po deadlinie nie ma potrzeby go odrabiać, jeśli wrócisz do zwykłego rytmu.

Awaria rytmu to inna sytuacja: jedna noc, w której nie zasypiasz do drugiej, a budzik dzwoni o szóstej. Deficyt wynosi $12 - 4 = 8$ godzin, czyli niewiele więcej niż przy deadlinie, ale nie jest rozłożony w czasie. Często nie chodzi o samą noc, tylko o to, co zjadło wieczór: ekran, rozmowa, zaległa praca. Zanim zaczniesz szukać rozwiązania w tabeli, zapisz, co się wydarzyło.

\begin{zasada}[Wracaj następnego wieczoru, nie nadrabiaj]
Po złej nocy wracasz do zwykłej pory zasypiania następnego wieczoru. Nie kładziesz się wcześniej, żeby odrobić stratę, i nie śpisz do południa. Jeden dzień nie jest kredytem, który trzeba spłacić w tej samej walucie.
\end{zasada}

Dlaczego nie nadrabiać? Nadrabianie przesuwa rytm. Po nocy z 4 godzinami snu, po której śpisz 14 godzin, następna noc zaczyna się później, a wyrównanie zaczyna się od przesunięcia pory, a nie od odpoczynku. Jeśli chcesz odrobić 8 godzin jedną nocą, dostajesz $12 + 8 = 20$ godzin w łóżku, czyli prawie cały dzień bez czynności. To nie jest liczba, którą plan przewiduje, i nie jest to rozwiązanie, które chcesz wdrożyć. \verify{jednorazowe wydłużenie snu nie przywraca pełnej sprawności po wcześniejszym deficycie, a regularność pory zasypiania bywa ważniejsza od samej sumy godzin}.

\begin{falle}[Weekendowe nadrabianie]
Najczęstszy błąd: od piątku do niedzieli śpisz trzy godziny dłużej, a w poniedziałek budzik wyrywa cię z innego rytmu. W tabeli tydzień wygląda dobrze, a w praktyce każdy poniedziałek zaczyna się od przesunięcia. Korekta: jeśli chcesz zmienić porę snu w weekend, zmieniaj ją o godzinę, nie o trzy, i trzymaj poranną godzinę wstawania bliżej dnia roboczego.
\end{falle}

Deadline, choroba i awaria mają wspólną cechę: każda z nich przesuwa tydzień, nie czyniąc z niego porażki. Plan nie musi być nietknięty. Musi dać się wznowić.

\section{Zakwas w czasie urlopu}

Zakwas to jedyny element planu, który żyje, zanim wypiekasz pierwszy chleb, i nie czeka na twój wolny tydzień. W przepisach z tabeli chleb pszenny na zakwasie używa 100 grama zakwasu, a żytnio-pszenny 150 gramów. Oba potrzebują zakwasu aktywnego, a nie takiego, który stał tydzień w szafce.

Przed urlopem zakwas idzie do lodówki po karmieniu. W lodówce jego aktywność zwalnia i można go na jakiś czas zostawić bez karmienia. \verify{zakwas przechowywany w lodówce może przetrwać kilka tygodni bez karmienia, ale przed pierwszym wypiekiem wymaga kilku karmień w temperaturze pokojowej}. Po powrocie nie wyjmuj mąki od razu. Najpierw wyjmij zakwas, nakarm go, poczekaj, aż zacznie pracować, i dopiero wtedy zaplanuj wypiek.

Tu liczy się zapas dni, nie godzin. Wypiek z przepisu zakłada cztery godziny fermentacji w cieple w przypadku chleba pszennego na zakwasie. Te cztery godziny zaczynają się dopiero wtedy, gdy zakwas pracuje, a to może zająć więcej czasu niż sam przepis. Zaplanuj więc pierwszy tydzień po powrocie jako tydzień z przygotowaniem, a nie z wypiekiem.

Jeżeli wyjeżdżasz często, zakwas nie musi być wyrzucany za każdym razem. Ale jeśli nie piekłeś od dłuższego czasu, a zakwas stał w lodówce bez karmienia, potraktuj go jak nowy: zrób dwa karmienia, zanim podejmiesz decyzję, czy go zachować.

\section{Przegląd miesięczny: pięć pytań}

Przegląd robisz raz w miesiącu, w ostatni dzień albo w pierwszy dzień następnego. Zajmuje dwadzieścia minut i opiera się na zapisie, nie na pamięci. Pamięć zniekształca: zapamiętujesz gorszą noc jako całą porażkę miesiąca albo dobry tydzień jako dowód, że wszystko działa. Zapis daje liczby, a z liczb da się podjąć decyzję.

Przegląd składa się z pięciu pytań, w tej kolejności:

\begin{enumerate}
\item Ile dni wypadło z planu i z jakiego powodu? Odpowiadasz słowami z zapisu: dziecko, podróż, choroba, deadline, awaria, nic.
\item Ile godzin snu wyszło średnio tygodniowo i w których tygodniach był największy rozrzut? Liczysz z tabeli, nie z wrażenia.
\item Który blok był najczęściej przesuwany? Praca, dom, chleb, ruch, czy coś, czego nie ma w tabeli.
\item Co kosztowało najwięcej energii i czy da się to usunąć, zmniejszyć albo przenieść na inną porę?
\item Jaką jedną zmianę wprowadzasz w następnym miesiącu? Jedną, nie pięć.
\end{enumerate}

Piąte pytanie jest najtrudniejsze, bo kusi, żeby odpowiedzieć listą. Lista pięciu zmian zostanie porzucona w drugim tygodniu i nie wiadomo będzie, która z nich zawiodła. Jedna zmiana da się ocenić na koniec następnego przeglądu.

\section{Kiedy zrezygnować z 12 godzin na rzecz 9}

Dwanaście godzin snu było wyborem, a wybór może się zmienić. Przejście na 9 godzin uwalnia $12 - 9 = 3$ godziny dziennie, czyli $3 \cdot 7 = 21$ godzin tygodniowo. Zobaczmy, co się dzieje z szablonem A z rozdziału 6. Po zmianie sen wynosi 9 godzin, praca 8, ruch i higiena 1, posiłki 1,5, dom 0,5, chleb 0,5. Suma to $9 + 8 + 1 + 1{,}5 + 0{,}5 + 0{,}5 = 20{,}5$ godziny. Brakuje $24 - 20{,}5 = 3{,}5$ godziny, i to właśnie trafia do bloku \emph{Wolne}, który wynosi wtedy 3,5 zamiast \val{week.A.free}. W szablonie B wolne rośnie z \val{week.B.free} do 3,25.

Kiedy warto to rozważyć? Nie wtedy, gdy masz gorszy tydzień. Warto wtedy, gdy przegląd pokazuje powtarzający się wzór: przez kilka miesięcy wypełniasz tabelę, ale sen wychodzi regularnie krótszy niż 12 godzin i nic nie wskazuje, żebyś po nim funkcjonował gorzej; albo 12 godzin wydłuża się w weekend, skraca w tygodniu, a rytm nie trzyma się przez żaden dłuższy okres; albo 12 godzin kosztuje wyraźnie coś, czego nie chcesz oddawać, czyli pracę, dziecko, zdrowie albo zwykłe wieczory z ludźmi.

Jest też druga strona. Jeśli po 12 godzinach nadal jesteś senny, zmniejszanie bloku nie rozwiąże problemu. Wtedy właściwym krokiem jest konsultacja z lekarzem, a nie kolejna korekta tabeli.

Zmiana z 12 na 9 nie jest porażką. Jest zmianą budżetu, a nie charakteru. Książka nazywa się od dwunastu godzin, bo taki był punkt wyjścia do rozmowy o dobie, a nie dlatego, że dwanaście godzin jest wartością samą w sobie. Część osób po przejściu na 9 godzin ma czas, którego wcześniej nie miała, i to może być lepszy wybór. Nie musisz zgadzać się z tym z góry. Sprawdź to w zapisie, przez miesiąc.

\section{Co zostaje z planu}

\val{chapters.total} rozdziałów to dużo słów o jednej prostej rzeczy: doba ma dwadzieścia cztery godziny i każdy blok jest zobowiązaniem. Jeśli z tej książki ma zostać jedna praktyka, to jest nią liczenie. Plan jest sumą nazwanych bloków, a nie obrazem idealnego dnia. Kiedy suma nie zgadza się z rzeczywistością, to rzeczywistość wygrywa, a plan się poprawia.

Zostaje kilka zasad, które przetrwają każdą zmianę. Wracasz następnego wieczoru zamiast nadrabiać. Zmieniasz jedną rzecz na raz. Przegląd opiera się na zapisie. Czas bierny, na przykład fermentacja, jest czasem, który można zaplanować, a nie tylko wypełnić. Dwanaście godzin snu to budżet, który możesz zmienić, a jego zmiana nie jest dowodem, że coś poszło źle.

Część liczb w tej książce policzysz sam, z własnej tabeli. Część pochodzi z przepisów i szablonów, które można sprawdzić. A część twierdzeń jest oznaczona jako do sprawdzenia, bo nie są one wyliczone w tekście i nie należy im wierzyć na słowo. Przy każdym z nich zapytaj, skąd się biorą, zanim wpiszesz je do swojego planu.

Plan nie musi się domknąć. Ma się dać wznowić, a to już wymaga nie tylko tabeli, ale też kilku nawyków, które pozwolą wrócić do rytmu po każdym tygodniu, który się nie ułożył.

\begin{cwiczenia}
\exitem{Trzy doby z niemowlęciem: sen 5, 6, 1 i 7 godzin, cel 12 godzin dziennie. Ile godzin brakuje?}{Cel to $3 \cdot 12 = 36$ godzin, wykonanie to $5 + 6 + 1 + 7 = 19$ godzin, brakuje $36 - 19 = 17$ godzin.}
\exitem{Szablon A po przejściu na 9 godzin snu. Ile godzin trafia do bloku Wolne?}{Suma bez wolnego to $9 + 8 + 1 + 1{,}5 + 0{,}5 + 0{,}5 = 20{,}5$ godziny, więc wolne wynosi $24 - 20{,}5 = 3{,}5$ godziny. Wcześniej było 0,5 godziny.}
\exitem{Trzy noce po 9 godzin snu zamiast 12. Ile godzin deficytu się zbiera i czy trzeba je odrobić jedną nocą?}{Deficyt to $3 \cdot 3 = 9$ godzin. Nie trzeba go odrabiać jedną nocą: odrobienie wymagałoby $12 + 9 = 21$ godzin w łóżku, a zasada mówi, żeby wrócić do zwykłej pory następnego wieczoru.}
\exitem{Chleb drożdżowy na zimno ma 12 godzin fermentacji w lodówce. Zaczynasz o 20:00. O której godzinie najwcześniej możesz zacząć pieczenie?}{$20 + 12 = 32$, czyli o 8:00 następnego dnia. Najwcześniej, bo fermentacja nie powinna być skracana, a pieczenie zaczyna się po niej.}
\exitem{Przelot sześć godzin na wschód. W miejscu docelowym jest 22:00. Która godzina jest w domu?}{$22 - 6 = 16$, czyli 16:00. Przeliczenie wykonujesz w jedną stronę, a nie w obie.}
\exitem{Dlaczego przegląd opiera się na zapisie, a nie na pamięci tygodnia?}{Pamięć zawyża albo zaniża jeden zły lub dobry tydzień i odtwarza go jako cały miesiąc. Zapis daje liczby, z których da się policzyć średnią i rozrzut.}
\exitem{Jeden miesiąc, pięć odpowiedzi w przeglądzie. Dlaczego na piąte pytanie odpowiadasz jedną zmianą, a nie listą?}{Lista pięciu zmian wprowadzonych naraz uniemożliwia ocenę, która z nich zadziałała albo zawiodła. Jedna zmiana daje się ocenić na następnym przeglądzie.}
\end{cwiczenia}

\zadanie{Wybierz w najbliższym miesiącu jeden tydzień, który wypadnie z planu: podróż, deadline, gorszy dzień w domu albo choroba. Przed nim zapisz, który blok zostanie przesunięty i ile godzin mu przypisujesz. Po tygodniu zsumuj godziny z zapisu, nie z pamięci, i porównaj wynik z budżetem tygodniowym 84 godziny snu. Miara sukcesu: każda doba w zapisie sumuje się do 24 godzin, a przegląd na koniec miesiąca ma pięć odpowiedzi, z których jedna prowadzi do konkretnej zmiany.}
